% !TEX root=jamc80.tex
\documentclass[document.tex]{subfiles}

\begin{document}
\section{\cite{jamc80} General equations}
Conservation of mass, momentum and energy: 

\textcolor{green}{ADD ALGEBRA COMPARISON SCHUBERT in final doc}
\begin{align}
    \pdv{\rho}{t} + \nabla \cdot (\rho\bm{u}) &= 0  \label{eq:3}\\
    \rho\left(\pdv{\bm{u}}{t} + \bm{u} \cdot \nabla \bm{u}\right) &= -\nabla P +\rho g + \pdv{x_j} \tau_{ij} \label{eq:4}\\
    \rho C_p \left\{\pdv{T}{t} - \frac{\alpha T}{\rho C_p}\pdv{P}{t} + \bm{u} \cdot [\nabla T - (\nabla T)_s]\right\} &= 
    \nabla \cdot (K\nabla T) + H + \tau_{ij}\pdv{u_i}{u_{ij}} \label{eq:5}
\end{align}

All terms and symbols used in these equations are shown in table \ref{table:symbols}. 
The (Newtonian) relation between the deviatoric stress tensor $\tau_{ij}$ and the velocity field is

\begin{equation}
\tau_{ij} = \eta \left( \frac{\partial u_i}{\partial x_j} + \frac{\partial u_j}{\partial x_i}
 - \frac{2}{3}\,\delta_{ij}\,\frac{\partial u_k}{\partial x_k} \right) 
 + \eta_2\,\delta_{ij}\,\frac{\partial u_k}{\partial x_k}.
\label{eq:8}
\end{equation}

Combining Eq.~\eqref{eq:8} with the momentum equation \eqref{eq:4} gives the general Navier-Stokes equations, 
but \textcite{jamc80} does not show these. Instead, they assume that $\frac{1}{Pr} \to 0$, which
they assume (\textcolor{green}{do they?}) is the case for the Earth's mantle, the inertial forces can be ignored. The resulting equations 
will be discussed in the next section, which covers the approximations made by \textcite{jamc80}. 
We also define the constant Grüneisen's parameter: 
\begin{align} \label{eq:10}
    \Gamma = \frac{\alpha K_s}{\rho C_p} = \frac{\alpha K_T}{\rho C_v}
\end{align}
Where $K_s$ and $K_T$ are the isentropic and isothermal moduli of bulk compressibility, 
and $C_v$ and $C_p$ are the specific heat at constant volume and pressure, respectively (\ref{table:symbols}).

\begin{mdframed}
    Additional equations used in derivations in the next section: 
    \begin{align} \label{eq:specific_heats}
        C_p - C_v = \frac{\alpha^2K_TT}{\rho}
    \end{align}
\end{mdframed}

\section{\cite{jamc80} Approximations}
Problem statement: for the Boussinesq approximation to be valid, the depth range of convection $d$ must be 
small compared to the temperature scale height $H_T$ ($d/H_T \ll 1$). This is not the case for many systems of 
deep convection, amongst which the convection of planetary mantles. Physically, $d/H_T \approx O(1)$ means that 
the influence of viscous dissipation, variable adiabatic temperature- and density gradients must be maintained in the 
set of governing equations. 
\newline
The assumptions made for the \textbf{anelastic liquid approximation} are that $d=O(H_T)$ and that $\alpha T \ll 1$
\textcolor{green}{what does this mean and how does this relate to Prandtl number? 
and if $\alpha \ll 1$, how come $D$ is not $\ll 1$?} 

\subsection{Approximate momentum equation}
Since volume changes in the Earth's mantle occur over timescales greater than $10^7$ years, terms containing the bulk viscosity 
$\eta_2$ are neglected. Substituting from Eq.~\eqref{eq:8}, the momentum equation, Eq.~\eqref{eq:4}, becomes
\begin{align} \label{eq:9}
    0 = -\nabla P + \rho \mathbf{g} + \pdv{x_i}\, \eta \left( \pdv{u_i}{x_j} + \pdv{u_j}{x_i} \right) - \frac{2}{3} \nabla \eta \, (\nabla \cdot \mathbf{u}).
\end{align}
\subsection{Reference state and equation of state}
A reference state with hydrostatic pressure $P_H$ and adiabatic temperature distribution (gradient?) $T_s$ is considered. 
The equation of state is written as a linear Taylor expansion of $\rho$ about the reference state density distribution $\rho_r$: 
\begin{equation} \label{eq:11}
    \rho(T,P) = \rho_r[1-\alpha(T-T_s)+K_T^{-1}(P-P_H)]
\end{equation}
Where 
\begin{align}
    \alpha &= \frac{1}{\rho_r} \left(\pdv{\rho_R}{T}\right)_P \\
    K_T &= \rho_r \left(\pdv{P}{\rho_r}\right)_T 
\end{align}

Determining the reference density distribution $\rho_r$ is done by considering the gradient of $\rho_r(S, P_H)$ with $S$ 
the entropy per unit mass: 
\begin{equation}
    \nabla\rho_R(S,P_H) = \left(\pdv{\rho_r}{S}\right)_{P_H} \nabla S + \left(\pdv{\rho_r}{P_H}\right)_S\nabla P_H
\end{equation}
With $\nabla S = 0$ since the reference state is adiabatic, and with $\nabla P_H = -\rho_rg \bm{\hat{z}}$:
\begin{equation}
    \frac{1}{\rho_r} \pdv{\rho_r}{z} = -\frac{\rho_r g}{K_s}
\end{equation} 
with $K_s = \rho_r\left(\pdv{P_H}{\rho_r}\right)_s$. Using the definition of the Grüneisen parameter $\Gamma$ and the temperature scale height of the liquid
$H_T$ gives the Adams-Williamson relation \textcolor{green}{add images!}
\begin{equation} \label{eq:adams-williamson}
    \frac{1}{\rho_r} \pdv{\rho_r}{z} = -\frac{1}{H_T \Gamma}
\end{equation}
Which is valid for a \textbf{chemically homogeneous adiabatic density distribution under hydrostatic pressure}. 
Integration of equation \ref{eq:adams-williamson} over $z$ gives the reference densitiy distribution as
\begin{equation} \label{eq:rho_r}
    \rho_r(z)= \rho_0 e^[(d-z)/H_T \Gamma]
\end{equation}
where $\rho_0 = \rho_r(z=d)$ is the reference density value at the surface \textcolor{green}{orientation of coordinate system??}
Plugging this expression of the reference density distribution into \ref{eq:11} gives the (linearized equation) of state
\begin{equation} \label{eq:16}
    \rho(T,P) = \rho_0 e^[(d-z)/H_T \Gamma][1-\alpha (T-T_s)+K_T^{-1}(P-P_H)]
\end{equation}

\subsection{Nondimensionalization} \label{subsection:nondimensionalization}
\ref{table:characteristic_values}, \ref{table:characteristic_values}. \textcolor{green}{move tables here?}
In this section, \cite{jamc80} introduces various characteristic values of variables in order to nondimensionalize the preceding equations. These characteristic values and the resulting dimensionless 
variables are presented in the tables below.The characteristic value for \textbf{velocity} is taken directly from \cite{mcrw74}, and their derivation is expanded below: 

\begin{mdframed} 
    Our starting point is the momentum equation (in the notation of \cite{mcrw74}), with the left-hand side set to 0 (on the basis of $\mathrm{Re} \ll 1$ and $\frac{1}{\mathrm{Pr}}\to 0$):
    \begin{align} \label{eq:mcrw:1}
        0 &= -\frac{1}{\rho} \pdv{P}{x_i} + \pdv{\phi}{x_i} + \frac{1}{\rho}\pdv{\tau_{ij}}{x_j}
    \end{align}
    with 
    \begin{align*}
        \tau_{ij} &= \rho \nu\left(\pdv{u_i}{x_j} + \pdv{u_j}{x_i}\right) \\
        \nabla\phi &= \bm{g}
    \end{align*}
    last term in Eq.~\eqref{eq:mcrw:1} becomes \textcolor{green}{help index confusion}
    \begin{align*}
        \frac{1}{\rho}\pdv{\tau_{ij}}{x_j} &= \rho \nu \left(\pdv{x_j}\pdv{u_i}{x_j}+\pdv{x_j}\pdv{u_j}{x_i}\right) \\
        &= \rho \nu (\nabla^2u_i+ \pdv{x_i}\nabla \cdot \bm{u}) \\
    \end{align*}
    \textbf{assuming incompressibility,} $\nabla \cdot \bm{u} = 0$, and \textcolor{green}{or should this be $\bm{u}$ (vector)?}
    \begin{align} 
        \frac{1}{\rho}\pdv{\tau_{ij}}{x_j} = \nu \nabla^2 u_i
    \end{align}

    Now we formulate an equation of state under the Boussinesq approximation (Taylor expansion of density $\rho(T)$ around a reference state temperature $T_0$): 
    \begin{align*}
        \rho &= \rho_0[1-\alpha (T-T_0)] \\
        \frac{1}{\rho} &= \rho_0^{-1}[1-\alpha (T-T_0)]^{-1} \\
        &\approx \rho_0^{-1}[1+\alpha (T-T_0)]
    \end{align*}
    with $\alpha = -\frac{1}{\rho}\left(\pdv{\rho}{T}\right)_P$ the thermal expansivity coefficient. 

    We can split the pressure gradient $\nabla P$ into a hydrostatic part $(\nabla P_s)$ and a non-hydrostatic part $\nabla p$: 
    \begin{align*}
        \nabla P &= \nabla P_s + \nabla p  \\
        &= \rho_0 \bm{g} + \nabla p
    \end{align*}
    Together with the equation of state, the pressure term + gravity in Eq.~\eqref{eq:mcrw:1} then becomes 
    \begin{align*}
        -\frac{1}{\rho} \nabla P + \bm{g}&= -\rho_0^{-1}(1+\alpha(T-T_0))(\nabla P_s + \nabla p) + \bm{g}\\
        & = -\rho_0^{-1} \nabla P_s  - \rho_0^{-1}\nabla P_s \alpha(T-T_0) - \rho_0^{-1} \nabla p -  \rho_0^{-1} \nabla p \alpha(T-T_0) + \bm{g} \\
        & \approx -\bm{g}\alpha(T-T_0) - \rho_0^{-1}\nabla p
    \end{align*}
    Where we neglected the second-order terms ($\sim \nabla p T$) and used that $\nabla P_s = \rho_0\bm{g}$, cancelling out the gravity terms.
    With this, Eq.~\eqref{eq:mcrw:1} becomes 
    \begin{align} \label{eq:mcrw:10}
        \nu \nabla^2 \bm{u} = \bm{g}\alpha(T-T_0)+\rho_0^{-1}\nabla p
    \end{align} 
    \textbf{note that} $\nabla p$ \textbf{is only the non-hydrostatic pressure gradient}. 

    We can take the curl of Eq.~\eqref{eq:mcrw:10}, using that $\nabla \times \bm{u} = \bm{\omega}$, that the curl of a gradient is always zero ($\nabla \times (\nabla f) =0$), and that $\bm{g} = \nabla \phi$: 
    \begin{align} \label{eq:mcrw:13}
        \begin{split}
            \nu \nabla^2\bm{\omega} &= \nabla \times (\bm{g}\alpha(T-T_0))+\nabla \times (\rho_0^{-1}\nabla p) \\
        &= \nabla \times [\bm{g}\alpha(T-T_0)] \\
        &= \alpha(T-T_0)(\nabla \times \bm{g})+ \alpha \nabla T \times \bm{g} \\
        & = \alpha \nabla T \times \bm{g}
        \end{split}
    \end{align}
Now we assume that the flow in restricted to the $x,z$ plane, meaning 
$\bm{u} = (u,0,w) = \left(-\pdv{\psi}{z}, 0, \pdv{\psi}{x}\right)$ (where $\psi$ is the stream function) 
and thus $\bm{\omega} = (0,\omega, 0) = (0, -\nabla^2\psi, 0)$. With $\bm{g} = (0,0,g)$, Eq.~\ref{eq:mcrw:13} becomes (computing the cross product)
\begin{align} \label{eq:mcrw:16}
    \nabla^2\omega = \frac{\alpha g}{\nu}\pdv{T}{x}
\end{align}
From this equation, we can get an characteristic velocity for the convecting system. 
\begin{align} \label{eq:U}
    \begin{split}
       \mathcal{O}(\nabla^2\omega) &=\mathcal{O}(\frac{\alpha g}{\nu} \pdv{T}{x}) \\
    \frac{U}{d^3} &= \frac{\alpha g}{\nu_0} \frac{T_1}{d} \\
    U &= \frac{\alpha g T_1 d^2}{\nu_0} 
    \end{split}
\end{align}
Where $U$, $\nu_0$, and $T_1$ are characteristic values for velocity, kinematic viscosity, and temperature deviations from the adiabatic reference profile $T_0$, respectively. 
\end{mdframed}

For the scaling of pressure variations $P_1$ are derived by considering the horizontal component of Eq.~\eqref{eq:9} and assuming the non-
hydrostatic pressure gradient $\nabla P$ and the viscous force terms to be of the same order of magnitude
\textcolor{green}{is this an illegal way to write this?}: 
\begin{align*}
    \mathcal{O}(\nabla P) &= \mathcal{O}\left(\pdv{x_i}\, \eta \left( \pdv{u_i}{x_j} + \pdv{u_j}{x_i} \right) - \frac{2}{3} \nabla \eta \, (\nabla \cdot \mathbf{u})\right) \\
    \frac{P_{1_r}}{d} &= \frac{\eta U}{d^2} \\
    P_{1_r} &= \frac{\eta U}{d}
\end{align*}
Thus, for the dimensionless pressure variable, using Eq.~\eqref{eq:U} and the kinematic viscosity $\nu_0 = \frac{\eta}{\rho_0}$, we get: 
\begin{align*}
    P_1' &= \frac{P_1}{P_{1_r}} = \frac{P_1 d}{\eta U} \\
    &= \frac{P_1}{\alpha T_1 \rho_0 gd}
\end{align*}

For the dimensionless energy generation rate $\epsilon_0 = \frac{H}{\rho C_p}$ we consider the heat generation and the time scale in the energy equation \eqref{eq:5}: 
\textcolor{green}{better way to write this down? order of magnitude way?}
\begin{align*}
    \pdv{T}{t} &\sim \epsilon_0 \\
    \frac{T_1}{\tau} \pdv{T'}{t'} &\sim \epsilon_0 \\
    \epsilon_0' &= \frac{\epsilon_0 \tau}{T_1} \\
    &= \frac{\epsilon_0 \nu_0}{T_1^2g\alpha d}
\end{align*}
Where expressions from table \ref{table:characteristic_values} were used to arrive at the final expression. 
Similarly for the thermal diffusivity, we balance the conduction term with the time-derivative term in the energy equation \eqref{eq:5} : 
\begin{align*}
    \mathcal{O}(\kappa_0 \nabla^2T) &= \mathcal{O}\left(\pdv{T}{t}\right) \\
    \kappa_{0,r} \frac{T_1}{d^2} &= \frac{T_1}{\tau} \\
    \kappa_{0,r} &= \frac{d^2}{\tau} \\
    \kappa_0' &= \frac{\kappa_0}{\kappa_{0,r}} \\
    &= \frac{\kappa_0 \nu_0}{d^3g\alpha T_1}
\end{align*}
Finally, for the dimensionless boxwidth, we simply divide the dimensional box width with the characteristic length: 
\begin{align*}
    \lambda = \frac{L}{d}
\end{align*}

\begin{table}[h!] 
    \renewcommand{\arraystretch}{1.8}   % taller rows
    \setlength{\tabcolsep}{12pt}        % wider column gaps
    \centering
    \caption{characteristic values of \cite{jamc80}}
    \label{table:nondimensionalisation}
    \begin{tabular}{l  >{$}l<{$} >{$}l<{$}}
        \toprule
        Variable & \text{Characteristic value} & \text{Dimensionless variable} \\
        \midrule
        depth & d & (x', z') = \dfrac{(x,z)}{d}\\
        temperature & T_1 (=1 \si{\kelvin}) & T' = \dfrac{(T-T_0)}{T_1}\\
        density & \rho_0 & \rho_r' = \dfrac{\rho_r}{\rho_0}, \quad \rho = \dfrac{\rho}{\rho_0} \\
        velocity & U = \dfrac{g\alpha T_1 d^2}{\nu _0} & \\
        kinematic viscosity & \nu_0 = \eta/ \rho_0 &  \\
        time & \tau = d/U & \\
        \bottomrule
    \end{tabular}
\end{table}




\subsection{The anelastic-liquid approximation}
From now on, only the dimensionless variables will be used but with primes surpressed \textbf{(i.e. all variables are assumed dimensionless)}. 
The dimensionless equation of state (from eq. \ref{eq:16}) becomes 
\begin{equation} \label{eq:19}
    \rho = \rho_r[1-\mu(T-T_s)+\mu(\rho_0gd/K_T)P_1]
\end{equation}
with $\mu = \alpha T_1$. If $\mu \to 0$, $\rho$ reduces to its reference profile $\rho_r$ (as is the case in the Boussinesq approximation). 

Now, we want to eliminate the bulk modulus from Eq.~\eqref{eq:19}. For this, we use the Grüneisen parameter \eqref{eq:10} and 
the specific heat identity \eqref{eq:specific_heats} to rewrite the last term: 
\begin{align*}
    \frac{\rho_0gd}{K_T} &= \frac{\rho_0gd}{\Gamma\rho C_p}\alpha(1+ \alpha\Gamma T) \\
    &= \frac{\rho_0}{\rho}\frac{d}{H_T\Gamma}(1+ \alpha\Gamma T)\\
    &= \frac{D}{\rho_r \Gamma}(1+\alpha\Gamma T)
\end{align*}
Where it was used that $\rho \to \rho_r$ and primes were dropped, and that $d/H_T = D$. 

Now we assume $T \approx T_s$ and \textcite{jamc80} (implicitly) introduces the adiabatic temperature profile: \textcolor{green}{NOT SURE ABOUT THIS}
\begin{align} \label{eq:adiabatic_temp}
    T_s = T_0e^{(d-z)/H_T} = T_0 e^{(1-z)D}
\end{align}
Where we used that $z' = z/d$ and dropped the prime. 
Subsituting \eqref{eq:adiabatic_temp} into \eqref{eq:19} gives \textcolor{green}{WHY WHY WHY derive this equation if you're just going to throw away all terms with $\mu$???}
\begin{align}
    \label{eq:20}
    \rho = \rho_r\left\{1-\mu(T-T_s) + \mu\left(\frac{D}{\rho_r \Gamma}\right)[1+\mu\Gamma T_0e^{(1-z)D}]P_1\right\}
\end{align}

\textbf{Now, some approximations are made:} 
\begin{itemize}
    \item First, the Mach number $M^2$ is evaluated. This dimensionless number measures the speed of the fluid relative to the local speed of sound. 
        Using the characteristic velocity $U$ derived in section \ref{subsection:nondimensionalization} (Table \ref{table:characteristic_values}) and 
        the velocity of sound $C = \sqrt{\frac{1.8K_s}{\rho}} = \sqrt{\frac{1.8C_p\Gamma}{\alpha}}$ \textcolor{green}{where does the 1.8 come from??}  
        \begin{align*}
        M^2 &= \frac{U^2}{C^2} \\
        &= \frac{d}{H_T} \alpha T_1 \frac{g\alpha T_1 d^3}{1.8\nu_0^2\Gamma} \\
        &= \frac{d}{H_T}\mu R* Pr^{-1} \frac{1}{1.8\Gamma}
        \end{align*}
        Now, for $\mathrm{O}(\Gamma) = 1$, $\mathrm{O}(\frac{d}{H_T})$, and $Pr \to \inf$, the Mach number $M^2 \ll 1$. 
        Based on this approximation, $\pdv{\rho}{t}$ is neglected form the mass continuity equation \eqref{eq:3}. \textcolor{green}{why? isentropic relation, work it out here?}
    \item Since $\alpha \approx 10^{-5} \si{\kelvin}^{-1}$ and $T_1 = 1 \si{\kelvin}$, $\mu \ll 1$ and from Eq.~\eqref{eq:20} we can see that 
            $\rho = \rho_r$ \textcolor{green}{best way to format units in math mode?}
\end{itemize}

Now, we can plug in the dimensionless $\rho_r$ \eqref{eq:rho_r} into the mass continuity equation \eqref{eq:3}: 
\begin{align}
    \begin{split}
        \nabla \cdot [\mathbf{u} e^((1-z)D/\Gamma)] &= 0 \\
        \nabla \cdot [\mathbf{u} e^(-zD/\Gamma)] &= 0 
    \end{split}
\end{align}
Since $\Gamma$ and $D$ are both constant. \textcolor{green}{check D to be constant to be sure. also why are they bold? and is D italic/plain?}

\textcolor{green}{so, for own sanity, a summary: the density in the mass continuity equation is reduced to the reference profile (neglecting variations away from density profile) while the full
density (including variations) are used in the momentum equation (where density is multiplied with gravity). In the boussinesq equation, is the density also reduced in the body force term? }

\subsection{Stream function and vorticity field}
\textcolor{green}{stream function for compressible flow? I thought it was only defined for incompressible flow? Or are we discussing 
incompressible flow here because the continuity equation is reduced, even though we have density variations in the other equations?}
\textbf{Restricting the flow to the $x$,$z$ plane} and using that $\mathbf{u} e^(-zD/\Gamma)$ is divergence-free (solenoidal), a stream function $\psi$ can be defined: 
\begin{align}
    \label{eq:23}
    \mathbf{u} e^{-zD/\Gamma} = (-\pdv{\psi}{z}, 0, \pdv{\psi}{x})
\end{align}
Expressing the velocity in terms of this stream function: 
\begin{align}
    \label{eq:24}
    \mathbf{u} = (u,0,w) = \left(-e^{zD/\Gamma}\pdv{\psi}{z},0,e^{zD/\Gamma}\pdv{\psi}{x}\right)
\end{align}
Now we can relate the stream function to the vorticity field $\mathbf{\omega}$, using that $\mathbf{\omega} = \nabla \times \mathbf{u}$: 
\begin{align*}
    \bm{\omega} &= -\left(\pdv{w}{x}-\pdv{u}{z}\right)\hat{y} = \omega \\
    \omega &=-e^{zD/\Gamma}\left[\pdv[2]{\psi}{x}+\pdv[2]{\psi}{z} + \frac{D}{\Gamma}\pdv{\psi}{z}\right] \\
    -\omega e^{-zD/\Gamma}&=\nabla^2\psi +\frac{D}{\Gamma} \pdv{\psi}{z} 
\end{align*}

Now, we can relate the vorticity to the temperature and the non-hydrostatic pressure by taking the curl of the momentum equation \eqref{eq:9} and substituting the 
(dimensionless) density distribution \eqref{eq:20}. This is a pretty lengthy derivation: 
\begin{mdframed}
    Starting point: Eq.~\eqref{eq:9}. Using that $\nabla \times (\nabla f) = 0$ and that $\nabla \times (\nabla \cdot \bm{A}) = 0$ for any $f$ or $\bm{A}$, taking the curl of Eq.~\eqref{eq:9} gives 
    \begin{align*}
        \nabla \times \left(-\nabla P + \rho \bm{g} + \pdv{x_i}\eta\left(\pdv{u_i}{x_j}+\pdv{u_j}{x_i}\right) - \frac{2}{3}\nabla \eta (\nabla \cdot \bm{u})\right)  = \nabla \times \left(\rho \bm{g}+\pdv{x_i}\eta\left(\pdv{u_i}{x_j}+\pdv{u_j}{x_i}\right) \right) \\
    \end{align*}
    Each term separately: 
    Since $\nabla \times \bm{g} = 0$ and we assume $\pdv{y} = 0$, 
    \begin{align*}
        \nabla \times (\rho\bm{g}) &= (\nabla \rho) \times \bm{g} \\
        &= g \pdv{\rho}{x}\hat{y}
    \end{align*}
    For the viscous term, we define: 
    \begin{align*}
        \pdv{x_i}\eta\left(\pdv{u_i}{x_j}+\pdv{u_j}{x_i}\right) = V_j
    \end{align*}
    Note that $j$ is a free index, while the summation is over $i=x,y,z$. Taking the curl thus gives
    \begin{align*}
        \nabla \times V &= -\left(\pdv{V_z}{x} -\pdv{V_x}{z}\right)\hat{y} \\
    \end{align*}
    Plugging in $V_x = 2\eta\pdv[2]{u}{x} + \eta \pdv{z}\left(\pdv{u}{z}+\pdv{w}{x}\right)$ (for which we carried out the summation over $i$. Note that since $\pdv{y} = 0$, the $i=y$ terms drop out.) and 
    $V_z = 2\eta\pdv[2]{w}{z} + \eta \pdv{x}\left(\pdv{u}{z}+\pdv{w}{x}\right)$ and simplifying: 
    \begin{align*}
        &= \eta\left[\pdv{z}(\pdv[2]{x}+\pdv[2]{z})u - \pdv{x}(\pdv[2]{x}+\pdv[2]{z})w\right] \\
        &=\eta \left(\pdv{z}\nabla^2u - \pdv{x}\nabla^2w\right)\\
        &= \eta\nabla^2\omega 
    \end{align*}
    Putting both terms together gives: 
    \begin{align*}
        \eta\nabla^2\omega = -g\pdv{\rho}{x}
    \end{align*}
    When nondimensionalized (see Table \ref{table:characteristic_values}), this becomes: 
    \begin{align*}
        \nabla^2\omega' &= -\frac{gd^2}{\nu_0U}\pdv[2]{\rho'}{x} \\
        &= -\frac{1}{\mu} \pdv{\rho'}{x'}
    \end{align*}
    Dropping the primes: 
    \begin{align*}
        \nabla^2\omega = \frac{1}{\mu}\pdv{\rho}{x}
    \end{align*}
    Now we can plug in (the dimensionless) Eq.~\eqref{eq:20} for $\rho$, where we use that $rho_r$, $T_s$, $\mu$, $D$, and $\Gamma$ are all constants: 
    \begin{align*}
            \nabla^2\omega &= -\frac{1}{\mu}\left[-\rho_r\mu\pdv{T}{x} + \mu \frac{D}{\Gamma}\left(1+\mu\Gamma T_0e^{(1-z)D}\right)\pdv{P_1}{x}\right] \\
            &= \rho_r\pdv{T}{x} -\frac{D}{\Gamma}\left(1+\mu\Gamma T_0e^{(1-z)D}\right)\pdv{P_1}{x}
    \end{align*}
    with $\mu \ll 1$ and $\rho_r = e^{(1-z)D/\Gamma}$ this reduces to 
    \begin{align} \label{eq:27}
        \nabla^2\omega = e^{(1-z)D/\Gamma}\pdv{T}{x}-\frac{D}{\Gamma}\pdv{P_1}{x}
    \end{align}
    We can determine $\pdv{P_1}{x}$ from the horizontal component of Eq.~\eqref{eq:9}: 
    \begin{align*}
        \pdv{P_1}{x} &= \eta \left(\nabla^2u+\frac{1}{3}\pdv{x}(\nabla \cdot \bm{u})\right) \\
    \end{align*} 
     In dimensionless form (with primes dropped): 
    \begin{align*}
        \pdv{P_1}{x} = \nabla^2u+\frac{1}{3}\pdv{x}(\nabla \cdot \bm{u})
    \end{align*}
    To eliminate $\nabla \cdot u$, we consider the continuity equation \eqref{eq:3} with $\pdv{\rho}{t} = 0$ and $\rho = \rho_r = e^{(1-z)D/\Gamma}$:
    \begin{align*}
        \nabla \cdot (\rho_r u) &= \rho_r\nabla \cdot u + u \cdot \nabla\rho_r \\
        &= \rho_r\nabla \cdot u - w \frac{D}{\Gamma} \rho_r = 0 \\
        \nabla \cdot u &= w\frac{D}{\Gamma}
    \end{align*} 
    Our (dimensionless) equation for $\pdv{P_1}{x}$ becomes
    \begin{align*}
        \pdv{P_1}{x} = \nabla^2u+\frac{1}{3}\pdv{x}\left(w\frac{D}{\Gamma}\right)
    \end{align*}
    Plugging this into Eq.~\eqref{eq:27} gives the vorticity equation: 
    \begin{align}
        \nabla^2\omega = e^{(1-z)D/\Gamma}\pdv{T}{x}-\frac{D}{\Gamma}\left(\nabla^2u+\frac{D}{3\Gamma}\pdv{w}{x}\right)
    \end{align}
\end{mdframed}




\section{Appendix 1: tables of symbols}
\begin{table}[ht] 
    \renewcommand{\arraystretch}{1.8}   % taller rows
    \setlength{\tabcolsep}{12pt}        % wider column gaps
    \centering
    \caption{Symbols used by \cite{jamc80}}
    \label{table:symbols}
    \begin{tabular}{>{$}l<{$} >{$}l<{$} l}
        \toprule
        d & & depth of the convecting layer \\
        H_T & \dfrac{C_p}{g\alpha} & temperature scale height of the liquid \\
        D & \dfrac{d}{H_T} & \\
        C_p & & specific heat at constant pressure \\
        C_v && specific heat at constant volume \\
        \alpha & & volume coefficient of thermal expansion \\
        \tau_{ij} && deviatoric stress tensor \\
        (\nabla T)_s & \dfrac{\alpha T}{\rho C_p}\nabla P & adiabatic temperature gradient \\
        K && thermal conductivity \\
        H && rate of internal heat production \\
        \eta & \rho \nu & dynamic viscosity \\
        \nu && kinematic viscosity \\
        \eta_0 & \ll 1 & second or bulk viscosity \\
        \Gamma & \dfrac{\alpha K_s}{\rho C_p} = \dfrac{\alpha K_T}{\rho C_v} & Grüneisen's parameter, assumed constant and $O(1)$\\ 
        K_s & & isentropic modulus of bulk compressibility \\
        K_T && isothermal modulus of bulk compressibility \\ 
        \mu & \alpha T_1 & \\
        \alpha & 10^{-5} \quad ^\circ K^{-1} & volume coefficient of thermal expansion \\ 




        % T_0 & & reference temperature \\
        % \Delta T^* & & temperature increase relative to reference temperature \\
        \bottomrule
    \end{tabular}
\end{table}

\begin{table}[ht] 
    \renewcommand{\arraystretch}{1.8}   % taller rows
    \setlength{\tabcolsep}{12pt}        % wider column gaps
    \centering
    \caption{characteristic values of \cite{jamc80}}
    \label{table:characteristic_values}
    \begin{tabular}{l  >{$}l<{$}  }
        \toprule
        depth & d\\
        temperature & T_1 = \dfrac{1}{\text{K}}\textcolor{green}{WHAT??} \to 1 ^{\circ} \text{K}\\
        density & \rho_0 \\
        velocity & U = \dfrac{g\alpha T_1 d^2}{\nu _0}, \quad \nu_0 = \eta/ \rho_0 \text{\textcolor{green}{help this is different from schubert}} \\
        time & \tau = d/U \\
        \bottomrule
    \end{tabular}
\end{table}

\begin{table}[ht] 
    \renewcommand{\arraystretch}{1.8}   % taller rows
    \setlength{\tabcolsep}{12pt}        % wider column gaps
    \centering
    \caption{Dimensionless variables and numbers in \cite{jamc80}}
    \label{table:dimensionless_variables}
    \begin{tabular}{>{$}l<{$} >{$}l<{$} }
        \toprule
        (x', z') & (x,z)/d\\
        \rho_r'& \rho_r/\rho_0 = e^[(1-z')D/\Gamma], \quad D= d/H_T\\
        \rho' & \rho/\rho_0 \\
        T'& (T-T_0)/T_1 \\ 
        t' & t/\tau \\
        P_1' & \dfrac{P_1 d}{U\eta} = P_1 \alpha T_1 \rho_0 gd, \quad P_1 = P-P_H \\
        \kappa'_0 & \kappa_0\nu_0/g\alpha T_1d^3, \quad \text{thermal diffusivity} \\
        \epsilon'_0 & \epsilon_0\nu_0/g\alpha T_1^2d, \quad \text{energy generation rate,}\quad \epsilon_0 = H/\rho_0C_p\\
        \lambda & L/d, \quad \text{box width, } \quad L \quad \text{is the dimensional box width}\\
        M^2 & U^2/C^2 \quad \text{Mach number, where $C$ is sound speed} \\
        R* & 1/\kappa_0' \quad \text{Rayleigh number equivalent} \\
        Pr & \nu_0'/\kappa_0' \quad \text{Prandtl number} \\
        \bottomrule
    \end{tabular}
\end{table}
\end{document}